% =============================================================================
% Section: Axiomatic Foundation
% =============================================================================

\section{Axiomatic Foundation}
\label{sec:axioms}

The \Opoch{} framework begins from $\noth$ (Nothingness: no committed
distinctions) and \emph{derives} its sole operating principle A0$^*$ as the
first theorem.  A0$^*$ is not declared---it is the unique admissibility
criterion forced by the absence of all external structure.

% ---------------------------------------------------------------------------
\subsection{Starting Point: $\noth$}
\label{sec:axioms:nothingness}

\begin{definition}[Nothingness $\noth$]
\label{def:nothingness-formal}
$\noth$ is the state with no committed distinctions:
no external labels, no external delimiter, no external clock, no external
witness, no primitive observer/observed split.  If anything else is assumed,
it is already not nothingness.
\end{definition}

% ---------------------------------------------------------------------------
\subsection{Derivation of A0$^*$ from $\noth$}
\label{sec:axioms:derivation}

From $\noth$, the admissibility criterion A0$^*$ is derived through a
chain of five necessity theorems:
\[
  \noth \;\to\; \text{N0} \;\to\; \text{N1} \;\to\; \text{N2}
  \;\to\; \text{N3} \;\to\; \text{N4} \;\to\; \text{A0}^*.
\]

\begin{theorem}[N0 --- No Externality]
\label{thm:N0}
At $\noth$, any proposed distinction that depends on an external
distinguisher is inadmissible, because no external distinguisher exists.
\begin{proof}
$\noth$ carries no committed distinctions.  An external distinguisher is
itself a distinction.  If it existed, $\noth$ would not be nothingness.
Contradiction.
\end{proof}
\end{theorem}

\begin{theorem}[N1 --- Endogenous Admissibility]
\label{thm:N1}
The only admissible distinction is one that can separate itself from
within reality---endogenously.
\begin{proof}
By N0, no external distinguisher exists.  If a distinction cannot be
separated from within, it requires an external agent to assert it, which
N0 forbids.  Therefore any admissible distinction must be endogenously
separable.
\end{proof}
\end{theorem}

\begin{theorem}[N2 --- Finiteness]
\label{thm:N2}
Any admissible witness must terminate in finite steps.
\begin{proof}
At $\noth$, there are no infinite resources---infinite resources would
constitute a committed external structure, violating $\noth$.  A witness
requiring infinite steps never produces a verdict, hence never separates
a distinction.  Therefore only finite witnesses are admissible.
\end{proof}
\end{theorem}

\begin{theorem}[N3 --- Replayability]
\label{thm:N3}
An admissible witness must carry its own replay trace.
\begin{proof}
Without an external recorder (forbidden by N0), the only way to verify a
witness's outcome is to replay it from the trace it produced.  A witness
that cannot be replayed depends on an unrepeatable external event,
violating N0.  Therefore replayability is forced.
\end{proof}
\end{theorem}

\begin{theorem}[N4 --- Internal Validity]
\label{thm:N4}
The distinction ``$w$ is a valid witness'' must itself be witnessable by
a finite endogenous procedure.
\begin{proof}
If witness validity were not witnessable, then ``$w$ is valid'' vs.\
``$w$ is not valid'' would be an external meta-judgment---a distinction
with no endogenous witness.  By N1, such a distinction is inadmissible.
Therefore witness validity must be endogenously checkable.
\end{proof}
\end{theorem}

\medskip\noindent
Theorems N0--N4 together force the unique admissibility criterion:

\begin{theorem}[A0$^*$ --- Completed Witnessability (Derived)]
\label{thm:a0star-derived}
A distinction is admissible if and only if a finite, endogenous,
replayable, self-accounting witness process whose own validity is
endogenously checkable can separate it.  A0$^*$ is the first theorem of
the framework, derived from $\noth$ via N0--N4---it is not postulated.
\end{theorem}

\begin{remark}[Set-Theoretic Operations as Instances of A0$^*$]
\label{rem:set-theory}
Basic set-theoretic operations (finite sets, membership, equality, union)
are themselves witnessable---each is a total computable function on finite
data---and are therefore members of the test set that A0$^*$ forces.  They are
not background assumptions; they are instances of the admissible procedures
that A0$^*$ demands.  The meta-language in which the paper is written (English,
mathematical notation) is not an axiom of the derivation, any more than
the choice of natural language is an axiom of a physics paper.
\end{remark}

\begin{remark}[Church--Turing as Historical Corroboration]
\label{rem:ct-status}
The Church--Turing thesis \citep{church1936,turing1936} independently
identified the class of total computable functions---the same class that
A0$^*$ forces at Step~3 (\cref{thm:executability}) via maximal closure under
totality-preserving composition.  The derivation does not use CT as a
premise; it derives the class directly as a mathematical theorem
(\cref{prop:utm-forced}).  The agreement between A0$^*$'s forcing and the
Church--Turing identification is historical corroboration, not a logical
dependency.  A reader who reads A0$^*$'s ``finite procedure'' as ``takes
finite input, runs finitely many steps, produces finite output'' arrives
at the same class without invoking CT.
\end{remark}

% ---------------------------------------------------------------------------
\subsection{From A0 to A0$^*$: The Closure Move}
\label{sec:axioms:closure-move}

The \emph{informal} starting point is:

\begin{quote}
\textbf{A0 (Bare Witnessability).}  A distinction is admissible if and only
if a finite witness procedure can separate it.
\end{quote}

This formulation is not yet self-consistent.  It speaks of ``valid witness
procedures'' without specifying what makes a witness valid.  The distinction
between ``$w$ is a valid witness for $\delta$'' and ``$w$ is not a valid
witness for $\delta$'' is itself a meaningful distinction.  By A0's own
criterion, that distinction must itself be witnessable by a finite
procedure internal to the universe.

Therefore an admissible witness process cannot be a bare separator only.
It must carry enough internal structure so that its own validity is finitely
checkable.  This is not a new axiom---it is A0 applied to its own witness
notion.  We call the result \textbf{A0$^*$ (Completed Witnessability)}.

% ---------------------------------------------------------------------------
\subsection{A0$^*$: Completed Witnessability (Derived)}
\label{sec:axioms:a0star}

\begin{remark}[A0$^*$ Is Derived, Not Postulated]
\label{rem:a0star-derived-status}
A0$^*$ is the first theorem of the framework, derived from $\noth$ via
the chain $\noth \to \text{N0} \to \text{N1} \to \text{N2} \to
\text{N3} \to \text{N4} \to \text{A0}^*$
(\cref{thm:N0,thm:N1,thm:N2,thm:N3,thm:N4,thm:a0star-derived}).
The \texttt{axiom} environment below is retained for referenceability,
but the content is a derived criterion, not a postulate.
\end{remark}

\begin{axiom}[A0$^*$ --- Completed Witnessability]
\label{ax:a0star}
A distinction $\delta$ is admissible if and only if there exists a finite
endogenous witness process $w$ satisfying all of the following:
\begin{enumerate}[nosep, label=(W\arabic*)]
  \item \textbf{Finite description:} $w$ has a finite endogenous code
        $\mathrm{code}(w) \in \Carrier$.
  \item \textbf{Replayability:} $w$ is replayable from its own trace---any
        agent with access to $\mathrm{tr}(w,s)$ can re-execute $w$ and
        obtain the same outcome.
  \item \textbf{Termination:} $w$ terminates within finite budget $B$
        on every admissible input.
  \item \textbf{Separation:} $w$ produces distinct outcomes on descriptions
        that $\delta$ distinguishes.
  \item \textbf{Disturbance accounting:} $w$ records its own state update
        $\mathrm{upd}(w,s)$---the change it induces on the state $s$---as
        part of the trace.
  \item \textbf{Endogenous validity:} the validity of $w$ is itself
        witnessable by a process satisfying (W1)--(W5) within the same
        universe.
  \item \textbf{Compositional closure:} if $w_1, w_2$ are admissible
        witness processes, then their ordered composition
        $w_1 \WitCompose w_2$ is again an admissible witness process.
  \item \textbf{Quotient invariance:} quotient-equivalent encodings of the
        same state induce equivalent witness traces---the witness does not
        depend on gauge-level representational choices.
\end{enumerate}
\end{axiom}

\begin{remark}[Bootstrap Reading of A0$^*$]
\label{rem:a0star-bootstrap}
The formal statement of A0$^*$ uses concepts (``carrier,'' ``code,''
``budget'') that are themselves derived in Steps~1--3.  This is not
circular: the \emph{informal} version---``a distinction is admissible iff
a finite, replayable, self-accounting witness process can separate it''---is
the primitive axiom.  The formal (W1)--(W8) listing is the \emph{derived}
restatement, expressed in the language that A0$^*$ itself forces.
\end{remark}

\begin{remark}[Why A0$^*$ and not bare A0]
\label{rem:why-a0star}
Bare A0 leaves the notion of ``valid witness'' as an external
meta-concept.  This creates a gap: the distinction between valid and
invalid witnesses would depend on an unwitnessed externality.  A0$^*$
closes this gap by demanding that witness validity be endogenously
checkable.  The additional structure (W1)--(W8) is not imported---it is
\emph{forced} by applying the witnessability criterion to itself.  Each
condition is individually necessary:
\begin{itemize}[nosep]
  \item Without (W1), witnesses are not finitely describable.
  \item Without (W2), witness outcomes are not independently verifiable.
  \item Without (W3), witnesses may consume infinite resources.
  \item Without (W4), witnesses do not actually separate.
  \item Without (W5), the distinction between ``state before $w$'' and
        ``state after $w$'' is unwitnessed.
  \item Without (W6), witness validity is an ungrounded meta-assertion.
  \item Without (W7), the test algebra is not maximally closed.
  \item Without (W8), truth depends on representational choices.
\end{itemize}
Omitting any one of (W1)--(W8) leaves a distinction that depends on an
unwitnessed externality, violating the witnessability criterion.
\end{remark}

A0$^*$ is the unique self-consistent completion of bare
witnessability, derived from $\noth$ via N0--N4.  It rules out distinctions that require infinite resources,
appeals to Platonic existence, unverifiable oracles, or untracked
disturbance.  In spirit it distills Bishop's constructivism
\citep{bishop1967} and Martin-L\"{o}f type theory \citep{martinlof1984}
into a single self-applied predicate, stated in computational rather than
type-theoretic terms.

% ---------------------------------------------------------------------------
\subsection{The Endogenous Boundary Theorem}
\label{sec:axioms:boundary}

\begin{theorem}[Endogenous Boundary --- Closed Universe is Forced]
\label{thm:endogenous-boundary}
Under A0$^*$, external delimiters, external verifiers, and external oracles
are inadmissible.  The universe must be informationally closed.
\begin{proof}[Proof sketch]
Suppose an external delimiter $v$ is used to parse witness traces.  Then
the distinction between ``same raw stream under delimiter $v_1$'' and
``same raw stream under delimiter $v_2$'' depends on $v$, which is not
itself an endogenous witness process satisfying (W1)--(W6).  This violates
(W6): the validity of any witness using $v$ is not endogenously
checkable.  The same argument applies to external verifiers and oracles.
Therefore all delimiters, verifiers, and composition rules must be
internally encoded and replayable.
\forcing{External dependencies create unwitnessed validity distinctions.
Endogenous closure is forced by (W6).}
\end{proof}
\end{theorem}

\begin{remark}[Closed Universe as L1, not L2]
\label{rem:closed-universe-l1}
In earlier drafts, the closed-universe assumption was classified as a
modeling commitment (L2).  \Cref{thm:endogenous-boundary} shows it is a
consequence of A0$^*$ (L1): the self-application of witnessability to
witness validity forces endogenous closure.  This eliminates one of the
former ``model-class commitments.''
\end{remark}

% ---------------------------------------------------------------------------
\subsection{The Ordered Witness Algebra}
\label{sec:axioms:witness-algebra}

A0$^*$ forces not just individual witness processes but an algebraic
structure on the collection of all admissible witnesses.

\begin{definition}[Ordered Witness Algebra $\WitAlg$]
\label{def:witness-algebra}
The \emph{ordered witness algebra} $\WitAlg$ is the set of all admissible
witness processes under A0$^*$, equipped with:
\begin{enumerate}[nosep, label=(\alph*)]
  \item \emph{Ordered composition:} $w_1 \WitCompose w_2$ denotes
        ``execute $w_1$, record its trace including state update, then
        execute $w_2$ on the updated state.''  This is associative but
        generally \emph{not} commutative.
  \item \emph{Involution:} $w^*$ denotes the \emph{replay adjoint}---the
        process that replays $w$'s trace in reverse, reading outcomes as
        inputs and inputs as outcomes.
  \item \emph{Positivity:} for every state $s$, the outcome weight
        $\GNSstate(w^* \WitCompose w) \geq 0$.
  \item \emph{Unit:} the identity witness $\mathbf{1}$ (do nothing, record
        nothing) is the composition identity.
  \item \emph{Norm:} $\|w\| = \sup_s \|w(s)\|$ is the operational
        supremum over admissible states.
\end{enumerate}
\end{definition}

\begin{proposition}[Witness Algebra is Forced]
\label{prop:witness-algebra-forced}
Each component of $\WitAlg$ is forced by A0$^*$:
\begin{itemize}[nosep]
  \item Ordered composition is forced by (W7) and (W5): composition must
        be closed, and state updates must be tracked, so execution order
        matters.
  \item The involution is forced by (W2): replayability requires the
        existence of a reverse-reading process for every witness.
  \item Positivity is forced by (W4): outcome weights from actual
        separations are non-negative.
  \item The unit is forced by (W7): the identity process is a trivially
        valid witness.
  \item The norm is forced by (W3): finite budget bounds all witness
        actions.
\end{itemize}
\forcing{The algebraic structure of $\WitAlg$ is the unique structure
compatible with (W1)--(W8).  Omitting any component leaves an algebraic
gap that permits unwitnessed distinctions.}
\end{proposition}

\begin{definition}[Witness Commutator]
\label{def:witness-commutator}
For witness processes $w_1, w_2 \in \WitAlg$, the \emph{commutator} is:
\[
  \WitCommutator{w_1}{w_2} \;=\;
  (w_1 \WitCompose w_2) - (w_2 \WitCompose w_1).
\]
When $\WitCommutator{w_1}{w_2} = 0$, the witnesses commute and their
execution order does not affect the truth quotient.  When
$\WitCommutator{w_1}{w_2} \neq 0$, the witnesses are \emph{non-commuting}
and order matters---each disturbs the state in a way that affects the
other's outcome.
\end{definition}

\begin{remark}[Noncommutativity is Internal from the Start]
\label{rem:noncommutativity}
The ordered witness algebra $\WitAlg$ is generally noncommutative because
(W5) tracks state disturbance.  Commutativity is not a global assumption
but a \emph{derived property} of specific witness pairs.  The commuting
sector $\CenterW$ (the center of $\WitAlg$) is the subalgebra of witnesses
that commute with all others.  The classical ``Diamond Law'' of earlier
drafts becomes a theorem about $\CenterW$, not a global axiom.  This
resolves the quantum gap: noncommuting observables are internal to the
framework from Step~0.
\end{remark}

% ---------------------------------------------------------------------------
\subsection{Consequences of A0$^*$}
\label{sec:axioms:consequences}

We now derive the structural consequences of A0$^*$ that form the substrate
for the full derivation in \cref{sec:derivation}.

% --- Carrier ---
\begin{definition}[Carrier Set $\Carrier$]
\label{def:carrier}
The \emph{carrier} is the set of all finite binary strings:
\[
  \Carrier \;=\; \{0,1\}^{<\infty}
    \;=\; \{\varepsilon,\; 0,\; 1,\; 00,\; 01,\; 10,\; 11,\; 000,\; \ldots\}.
\]
\end{definition}

\begin{proposition}[Carrier is Forced]
\label{prop:carrier-forced}
Under A0$^*$, every admissible witness has finite length (W1, W3).
Therefore every admissible description is a finite string over a finite
alphabet.  Up to recoding, this is $\Carrier = \{0,1\}^{<\infty}$.

\forcing{Any finite alphabet $\Sigma$ with $|\Sigma| \geq 2$ is
inter-translatable with $\{0,1\}$ by a computable bijection (W8).  An
alphabet of size~1 admits no distinctions.  Binary is therefore the
minimal non-trivial choice, and all others reduce to it.}
\end{proposition}

% --- Self-delimitation ---
\begin{proposition}[Self-Delimitation is Forced by Replayability]
\label{prop:prefix-forced}
A witness trace must be replayable from its own finite record (W2).  If
boundaries are not internally decidable, then the distinction between
``trace $= x$'' and ``trace begins with $x$'' is not witnessable from the
trace itself.  Therefore self-delimiting (prefix-free) coding is forced.
Any valid encoding must satisfy the Kraft inequality \citep{kraft1949}:
\[
  \sum_{x \in S} 2^{-|x|} \;\leq\; 1.
\]
\forcing{Self-delimitation is not a ``closed-universe extra.''  It is A0$^*$
applied to replayability (W2).}
\end{proposition}

% --- Endogenous tests and executability ---
\begin{definition}[Endogenous Witness Instruments]
\label{def:endogenous-tests}
An \emph{endogenous witness instrument} is a witness process
$w \in \WitAlg$ with:
\begin{itemize}[nosep]
  \item a code $\mathrm{code}(w) \in \Carrier$,
  \item an outcome map $\mathrm{out}(w, s) \in \Carrier$ for each state $s$,
  \item a state-update map $\mathrm{upd}(w, s)$ recording disturbance,
  \item a cost $c(w) \in \mathbb{N}_{>0}$.
\end{itemize}
When the update map is trivial ($\mathrm{upd}(w,s) = s$ for all $s$), we
call $w$ a \emph{pure test} (non-disturbing).  The set of all pure tests
is denoted $\Del$.  The full set of witness instruments (including
disturbing ones) is $\WitAlg$.
\end{definition}

\begin{definition}[Executability Primitives]
\label{def:executability}
The \emph{executability substrate} consists of:
\begin{enumerate}[nosep, label=(\alph*)]
  \item A \emph{total evaluator} $\UTM$ such that $\UTM(w, s)$ halts
        for all $w, s \in \Carrier$, returning both outcome and state
        update.
  \item A \emph{total decoder} $\mathrm{dec} \colon \Carrier \to
        \Carrier \times \Carrier$ parsing self-delimiting streams.
  \item A \emph{cost functional} $c \colon \WitAlg \to \mathbb{N}_{>0}$
        with $c(w) \geq 1$.
\end{enumerate}
\end{definition}

\begin{proposition}[Executability is Forced]
\label{prop:utm-forced}
A0$^*$ demands that \emph{all} finite witness procedures be admissible (W7).
The closure of all such procedures under composition and finite description
is precisely a universal evaluator $\UTM$ \citep{rogers1967}.

\forcing{If some finite procedure were excluded from the admissible set,
then there would exist a distinction (between ``inside'' and ``outside''
the admissible set) that no admissible witness could detect---violating
A0$^*$.  Hence the set of witnesses must be maximally closed.  Maximal
closure under totality-preserving composition on a countable carrier
yields exactly the total computable functions, and a universal evaluator
is forced.}
\end{proposition}

\begin{definition}[Finite Slice $\FinSlice$]
\label{def:finite-slice}
For a given budget $B$, the \emph{finite slice} is the operationally
reachable subset of the carrier:
\[
  \FinSlice \;=\; \bigl\{\, s \in \{0,1\}^* \;\bigm|\; |s| \leq B \,\bigr\}
    \;\subset\; \Carrier.
\]
All truth-quotient operations and witness executions are restricted to
$\FinSlice$.  The full carrier $\Carrier$ serves as the ambient space.
\end{definition}

\begin{remark}[The Church--Turing Corroboration]
\label{rem:ct-dependency}
\Cref{prop:utm-forced} derives the class of total computable functions
from A0$^*$ alone: maximal closure under totality-preserving composition on
a countable carrier is a mathematical theorem, not an empirical claim
\citep{rogers1967}.  The Church--Turing thesis independently identified
this same class by a different route---confirming the derivation's
conclusion, not entering it as a premise.
\end{remark}

% ---------------------------------------------------------------------------
\subsection{Summary of the Axiomatic Base}
\label{sec:axioms:summary}

The framework has \emph{no declared axioms}---A0$^*$ is the first
theorem, derived from $\noth$ via the necessity chain N0--N4.
\Cref{tab:axiom-summary} collects the logical dependencies.  Everything
downstream of this section is \emph{derived}, not postulated.  Notably,
all items formerly classified as ``modeling commitments'' (L2) are now
derived consequences of A0$^*$ (L1):

\begin{itemize}[nosep]
  \item \emph{Closed universe:} forced by (W6) via
        \cref{thm:endogenous-boundary}.
  \item \emph{Self-delimitation:} forced by (W2) via replayability.
  \item \emph{Ordered witness ledger:} forced by (W5) and (W7) via
        disturbance accounting and compositional closure.
  \item \emph{Noncommutativity:} internal from (W5); commutativity is a
        derived sector property.
\end{itemize}

\begin{table}[ht]
\centering
\caption{Axiomatic base and immediate consequences under A0$^*$.}
\label{tab:axiom-summary}
\begin{tabular}{@{}lll@{}}
\toprule
\textbf{Item} & \textbf{Status} & \textbf{Forced by} \\
\midrule
A0$^*$ (Completed Witnessability) & \emph{Derived from $\noth$} & N0--N4 \\
\midrule
Carrier $\Carrier = \{0,1\}^{<\infty}$ & \emph{Derived} & (W1), (W3) \\
Finite slice $\FinSlice \subset \Carrier$ & \emph{Derived} & (W3), Budget $B$ \\
Self-delimiting syntax $\SdMap$ & \emph{Derived} & (W2) \\
Closed universe / endogenous boundary & \emph{Derived} & (W6) \\
Ordered witness algebra $\WitAlg$ & \emph{Derived} & (W5), (W7) \\
Witness instruments (incl.\ disturbing) & \emph{Derived} & (W4), (W5) \\
Pure tests $\Del \subset \WitAlg$ & \emph{Derived} & (W4), $\mathrm{upd} = \mathrm{id}$ \\
Executability ($\UTM$, dec, $c$) & \emph{Derived} & (W1), (W3), (W7) \\
Noncommutativity (internal) & \emph{Derived} & (W5) \\
\bottomrule
\end{tabular}
\end{table}

With these objects in hand---a carrier $\Carrier$ with finite slice
$\FinSlice$, a self-delimiting syntax $\SdMap$, an ordered witness algebra
$\WitAlg$ with pure-test subalgebra $\Del$, executability primitives
($\UTM$, decoder, cost functional), and endogenous closure---we have the
complete substrate from which the derivation of \cref{sec:derivation}
proceeds.
